% !PRINTED 92
\weberpage{92}
A third kind of development, in which the symmetry of the
functions with respect to the two variables $v$, $a$ finds
expression, arises if the fractions that occurred in the previous
developments are developed in powers of $q$. In this way one
obtains
\[
\begin{gathered}
\frac{q^m}{e^{-2\pi i v}-q^m} = \Sigma\, q^{m\nu} e^{2\pi i \nu v}, \\
\frac{q^m}{e^{+2\pi i v}-q^m} = \Sigma\, q^{m\nu} e^{-2\pi i \nu v}, \\
\nu = 1,\, 2,\, 3,\, \ldots
\end{gathered}
\]
and these developments hold as long as the absolute value
of $q\, e^{2\pi i v}$ is a proper fraction, hence, if we put $v = v' + i v''$,
as long as
\[
-\omega'' < v < \omega'' % [sic] German source prints v; the imaginary-part bound mathematically requires v''.
\]
holds. Accordingly there follows from (14)
\[
\begin{gathered}
\frac{\vartheta'_{11}\,\vartheta_{11}(v+a)}{\pi\,\vartheta_{11}(v)\,\vartheta_{11}(a)} = \mathrm{cotg}\,\pi v + \mathrm{cotg}\,\pi a \\
{}+ 4\Sigma\, q^{\frac{mm'}{2}} \sin(m a + m' v)\pi,
\end{gathered}
\]
in which $m$, $m'$ run, independently of one another, through the \gesperrt{even} numbers
2, 4, 6, 8, \ldots\ This formula is valid for \gesperrt{real}
$a$ and $v$, and it holds beyond that as well, as long as the imaginary
part, of $v$ as well as of $a$, is in absolute value less than $\omega''$.

In Table III the 16 formulae of this kind are collected\textsuperscript{1}).

\textsuperscript{1}) Such developments have been considered by \gesperrt{Jacobi} (sur la rotation d'un corps)
(ges. Werke Bd.\ II), thereupon by \gesperrt{Hermite} (Annales de l'école normale 1885)
and by \gesperrt{Kronecker} (Berliner Akademie 1885). Cf.\ also the
Strasbourg dissertation of L.\ \gesperrt{Vockerodt} (1905).
